%% kapa: exam=1 round=27 year=2016 period=1 form=A subject=real_estate subject_ko=부동산학원론 seq=39 number=119
%% source: Q-Net 감정평가사 기출문제 (cst00309) artlSeq=5169402 file=extracted/1008/5169402/제27회 감정평가사 1차 문제/제27회 감정평가사 1교시문제(A형).pdf
%% conversion: kordoc 4.24.51 -> tools/gen1.py (auto); review=proofread-sonnet; flags=none
\begin{question}{119}
도시성장구조이론에 관한 설명으로 옳지 않은 것은?
\begin{choices}
  \item 버제스(Burgess)의 동심원이론은 도시생태학적 관점에서 접근하였다.
  \item 해리스(Harris)와 울만(Ullman)의 다핵심이론은 도시가 그 도시내에서도 수개의 핵심이 형성되면서 성장한다는 이론이다.
  \item 동심원이론은 도시가 그 중심에서 동심원상으로 확대되어 분화되면서 성장한다는 이론이다.
  \item 다핵심이론과 호이트(Hoyt)의 선형이론의 한계를 극복하기 위해서 개발된 동심원이론에서 점이지대는 저소득지대와 통근자지대 사이에 위치하고 있다.
  \item 선형이론은 도시가 교통망을 따라 확장되어 부채꼴 모양으로 성장한다는 이론이다.
\end{choices}
\end{question}
\answer{4}
